\ifx\allfiles\undefined% 如果这个没定义主文件
    \documentclass[12pt,a4paper,oneside]{ctexbook}
    \def\basepath{../}%basepath =  '../config'
    \input{../config/config}% 传进去后有 ../config/package.tex
    \begin{document}
\else% 如果定义了
\fi

\section{不等式}

本节介绍常用的一些不等式。

\medskip
\noindent\textbf{涉及绝对值的不等式}

实数 $a$ 的绝对值 $|a|$ 定义为
\[
|a| =
\begin{cases}
a, & \text{如果 } a\ge 0,\\
-a, & \text{如果 } a<0.
\end{cases}
\]

我们来考察不等式
\[
|x|<\alpha.
\]
根据绝对值的定义，该不等式等价于
\[
0\le x<\alpha,\quad \text{或者}\quad x<0,\ -x<\alpha,
\]
即
\[
-\alpha<x<\alpha.
\]
我们得到：
\[
|x|<\alpha \Longleftrightarrow -\alpha<x<\alpha.
\]
即：不等式 $|x|<\alpha$ 的解的集合是开区间
\[
(-\alpha,\alpha).
\]

类似地可以证明
\[
|y|\le \beta \Longleftrightarrow -\beta\le y\le \beta.
\]
即：不等式 $|y|\le \beta$ 的解的集合是闭区间
\[
[-\beta,\beta].
\]

我们有显然的不等式
\[
-|a|\le a\le |a|,\qquad -|b|\le b\le |b|.
\]
将这些不等式相加可得
\[
-(|a|+|b|)\le a+b\le |a|+|b|.
\]
由此又得到重要的不等式
\[
|a+b|\le |a|+|b|.
\]

运用这样的不等式，可以得到
\[
|a|=|(a-b)+b|\le |a-b|+|b|,
\]
\[
|a|-|b|\le |a-b|.
\]
同理可得
\[
|b|-|a|\le |b-a|=|a-b|,
\]
即
\[
|a|-|b|\ge -|a-b|.
\]
我们得到了
\[
-|a-b|\le |a|-|b|\le |a-b|,
\]
即
\[
\bigl||a|-|b|\bigr|\le |a-b|.
\]

利用归纳法，可以把不等式
\[
|a+b|\le |a|+|b|
\]
推广到 $n$ 个实数的情形：
\[
|a_1+a_2+\cdots+a_n|\le |a_1|+|a_2|+\cdots+|a_n|.
\]

\medskip
\noindent\textbf{伯努利(Bernoulli)不等式}

设 $x\ge 0$，则由二项式定理
\[
(1+x)^n
=1+nx+\frac{n(n-1)}{2}x^2+\cdots+x^n
\]
可以得到
\[
(1+x)^n\ge 1+nx.
\]
该不等式实际上对任何 $x\ge -1$ 成立。请看下面的定理。

\medskip
\textbf{定理}\quad 以下的伯努利不等式成立：
\[
(1+x)^n\ge 1+nx,\qquad \forall x\ge -1.
\]

\textbf{证明}\quad
我们采用数学归纳法。$n=1$ 时，上式显然以等式的形式成立。
假设已经证明了
\[
(1+x)^{n-1}\ge 1+(n-1)x,\qquad \forall x\ge -1,
\]
则
\begin{align*}
(1+x)^n
&=(1+x)^{n-1}(1+x)\\
&\ge [1+(n-1)x](1+x)\\
&=1+(n-1)x+x+(n-1)x^2\\
&\ge 1+nx,\qquad \forall x\ge -1.
\end{align*}
这证明了对一切自然数 $n$ 伯努利不等式成立。\qed

\medskip
\noindent\textbf{算术平均数与几何平均数不等式}

设 $x_1$ 和 $x_2$ 是非负实数。我们把
\[
\frac{x_1+x_2}{2}
\]
叫作这两个数的算术平均数，把 $\sqrt{x_1x_2}$ 叫作这两个数的几何平均数。
以下的不等式显然成立：
\[
(\sqrt{x_1}-\sqrt{x_2})^2\ge 0.
\]
由此得到
\[
\frac{x_1+x_2}{2}\ge \sqrt{x_1x_2}.
\]
即：算术平均数大于或等于几何平均数。

对 $n$ 个非负实数，也有相应的结果：

\medskip
\textbf{定理}\quad
设 $x_1,x_2,\cdots,x_n\ge 0$，则以下的算术平均数与几何平均数不等式
（AM-GM 不等式）成立：
\[
\frac{x_1+x_2+\cdots+x_n}{n}
\ge \sqrt[n]{x_1x_2\cdots x_n}.
\]

\textbf{证明}\quad
我们利用数学归纳法。$n=1$ 的时候，上式显然以等式形式成立。
假设对任意 $n-1$ 个非负实数，算术平均数与几何平均数不等式成立。
我们来考虑 $n$ 个非负实数 $x_1,x_2,\cdots,x_n$ 的情形。
不妨设 $x_n$ 是这 $n$ 个数中最大的一个
（我们总可以从小到大排列这 $n$ 个非负实数，于是最后一个数就是最大的一个）。
记
\[
A=\frac{x_1+\cdots+x_{n-1}}{n-1},
\]
则有
\[
x_n\ge A
=\frac{x_1+\cdots+x_{n-1}}{n-1}
\ge \sqrt[n-1]{x_1\cdots x_{n-1}}.
\]
于是
\begin{align*}
\left(\frac{x_1+\cdots+x_n}{n}\right)^n
&=\left[\frac{(n-1)A+x_n}{n}\right]^n\\
&=\left(A+\frac{x_n-A}{n}\right)^n\\
&=A^n+nA^{n-1}\left(\frac{x_n-A}{n}\right)+\cdots\\
&\ge A^n+nA^{n-1}\left(\frac{x_n-A}{n}\right)\\
&=A^n+A^{n-1}(x_n-A)\\
&=A^{n-1}x_n\\
&\ge x_1\cdots x_{n-1}x_n.
\end{align*}
即
\[
\frac{x_1+\cdots+x_n}{n}\ge \sqrt[n]{x_1\cdots x_n}.
\]
\qed

\medskip
\noindent\textbf{涉及三角函数的不等式}

几何学的讨论也是发现和证明不等式的一个重要途径。
我们举以下的例子来说明这种方法。在数学理论的推导中，
涉及角的度量时，通常采用弧度作为单位。

\medskip
\textbf{定理}\quad
对于用弧度表示的角 $x$，有以下不等式成立
\[
\sin x<x<\tan x,\qquad \forall x\in \left(0,\frac{\pi}{2}\right).
\]

\textbf{证明}\quad
在单位圆 $O$ 中作圆心角 $x$，它的始边为 $OX$ 轴上的 $OA$，
终边为 $OB$。用线段联结 $AB$。过 $A$ 点作 $OX$ 轴的垂线交 $OB$ 延长线于 $C$（图 1-1）。

\begin{center}
    \centering
    \includegraphics[width=0.42\textwidth]{1_5.jpg}

    {\small 图 1-1}
\end{center}

我们有：
\[
\triangle OAB \text{ 的面积}
< \text{扇形 } OAB \text{ 的面积}
< \triangle OAC \text{ 的面积},
\]
即
\[
\frac{1}{2}\sin x<\frac{1}{2}x<\frac{1}{2}\tan x.
\]
这样，我们证明了
\[
\sin x<x<\tan x,\qquad \forall x\in \left(0,\frac{\pi}{2}\right).
\]
\qed

\medskip
\textbf{推论}\quad
以下不等式成立
\[
|\sin x|\le |x|,\qquad \forall x\in\mathbb{R}.
\]

\textbf{证明}\quad
我们有
\[
|\sin x|=\sin x\le x=|x|,
\qquad \forall x\in\left[0,\frac{\pi}{2}\right),
\]
和
\[
|\sin(-x)|=\sin x\le x=|-x|,
\qquad \forall x\in\left[0,\frac{\pi}{2}\right).
\]
综合以上两式就得到
\[
|\sin x|\le |x|,
\qquad \forall x\in\left(-\frac{\pi}{2},\frac{\pi}{2}\right).
\]
而当 $|x|\ge \frac{\pi}{2}$ 时，又有
\[
|\sin x|\le 1<\frac{\pi}{2}\le |x|.
\]
于是，我们得到不等式
\[
|\sin x|\le |x|,\qquad \forall x\in\mathbb{R}.
\]
\qed

\ifx\allfiles\undefined %若没定义主文件
\end{document}% 就打印这个结尾
\fi